\chapter{Introduction}

\section{Goal of the Thesis}

The goal of this thesis is to develop and evaluate a reproducible quantization and hardware validation workflow for a Mamba-based neural network. The workflow connects four layers:

\begin{itemize}
  \item Python model export and hardware asset generation.
  \item C++ reference inference for deployment-oriented quantization.
  \item Python hardware-like simulation for RTL-oriented fixed-point behavior.
  \item RTL testbench validation using generated memory assets and golden outputs.
\end{itemize}

The main research questions are:

\begin{enumerate}
  \item Which parts of the model are sensitive to reduced precision?
  \item How can C++ reference results be used to separate quantization loss from hardware implementation loss?
  \item How can stage-level error analysis identify the dominant fixed-point error source?
  \item How does scaled-state quantization improve the selective scan path?
\end{enumerate}

\section{Contributions}

The main contributions are:

\begin{itemize}
  \item A Python--C++--RTL quantization flow with reproducible case assets.
  \item A C++ reference sweep for identifying int8-sensitive and int16-safe modules.
  \item A stage-level comparison method using \texttt{cpp\_like} and \texttt{hw\_like} traces.
  \item A scaled-state selective scan representation that reduces the dominant scan-stage mean absolute error.
\end{itemize}

\section{Thesis Outline}

\Cref{ch:background} introduces neural network quantization, fixed-point arithmetic, and the Mamba block structure. \Cref{ch:method} presents the proposed Python--C++--RTL workflow and the error metrics. \Cref{ch:results} discusses the quantization sweep and the selective scan scale optimization. \Cref{ch:conclusion} concludes the thesis and outlines future work.

